\justifying
\NSFCBodyText

\subsubsection{年度研究计划}
本课题按实验室开放课题组织，研究期限为\textbf{【暂定 24个月，自立项批复确定的起始时间计算】}。指南允许1—2年，实际起止日期和中期报告节点按批复执行。

\textbf{第一年度：完成主数据、失效诊断和核心机制验证。}
第1—3个月验收数据归档及矿区元信息，审计官方划分与权重覆盖，在主协议训练及验证矿区复现简单孪生与官方方法配置，最终测试保持封存。形成配准、标签边界及误报漏检清单，核对背景外观混淆是否具有足够普遍性；据此冻结数据划分、主要对照及指标定义。若数据或基线尚不支持机制假说，先处理质量问题并收窄研究目标。

第4—9个月完成同骨干四组约束对照、等面积匹配及平滑诊断，复现关键近邻方法，开展小规模掩码污染与特征行为分析。各配置保留随机种子、参数和日志，按矿区汇总效果。同步形成论文的问题、方法和实验初稿，避免把写作全部推迟到结题阶段。

第10—12个月根据开发结果冻结检测器，完成独立校准与错误排序原型，复核指标分母及数据职责。核心机制结果具备独立成稿条件时，争取第12个月完成首投；尚不具备时据实记录原因，以第15个月作为重新核对首投条件的节点。第一年度形成可复现实验版本、阶段性误差证据和核心消融，按实验室要求准备中期材料；阶段报告如实包含无增量结果与尚未解决的条件。

\textbf{第二年度：完成可靠性评价、统计分析与成果发表。}
第13—15个月完成冻结协议下的矿区评估、固定面积复核模拟和置信区间分析，集中检查高置信度漏检、细小变化及困难矿区。主链路闭合后，再根据资源决定EuroMineNet内部验证及专业人员案例复核；二者均不替代主数据核心对照。

第16—18个月集中处理审稿意见或调整投稿安排，第19—24个月继续修改并争取完成符合指南要求的发表，整理代码、划分、环境和统计记录，提交结题报告。论文评审周期存在不确定性，因此前期写作与实验同步推进，并保留必要的修改时间。成果署名、基金名称和批准号标注按正式管理要求落实。

\subsubsection{预期研究结果}
预期形成一致性约束对矿山变化判别的适用条件与反例，给出背景误报、小变化召回及约束分配之间的可检验关系；形成包含校准、错误排序与全图剩余错误的复核评价结果。效应以明确基线、有效矿区数及区间报告，不预设必须出现正向提升。

按指南要求发表至少1篇SCI/EI源刊论文，提交中期进展及结题报告；仅投稿、录用或EI会议论文不能自动作为已满足源刊发表要求的证明。附加成果为一套可复现的实验与评测流程，以及具有来源和解释边界的案例清单。代码和衍生产物公开前核对数据许可、既有项目权属和实验室成果共享约定。项目不将矿企现场部署、环境风险分级系统或新疆示范区建设列为必达成果。
